\section{Cartier trace and controlled lifting}
\label{c:trace-section}

The strict inequality from the previous section supplies a nonzero trace
on a geometric generic fibre of a full initial field. To turn it into
the section estimates needed later, we must lift that nonvanishing twice
and keep its exact value formula. This section proves each lifting step,
then checks that the finite spaces used to count all possible outputs
can be retained before reduction.

We write $\mathfrak p$ for the residue characteristic, reserving $p$
for the auxiliary cone dimension used earlier.  We write $P_0=\mathfrak p^l$ for a Frobenius
power.  All valuations, optimizers, strict discrepancy inequalities and
finite output bounds are chosen in characteristic zero.  Reduction transports
the specified finite algebraic data and numerical monomial formulas; it does
not assert the existence of an optimizer with the same properties in positive
characteristic.

The applications seek nonzero traced outputs $F'$ from actual inputs
$G\in V_M$.  In the retained reductions, survival of the $w$-initial
trace will give
\[
 F'\in V_{m'},\qquad m'\le M/P_0+n,\qquad
 w(F')=\frac{w(G)}{P_0}
       -\left(1-\frac1{P_0}\right)A_\theta(w).
\]
The corresponding lower value inequality will remain valid when the
positive weights vary in the same open stratum cone.  Thus inputs whose
degrees and values are $O(P_0)$ give outputs in degree and value ranges
bounded independently of $P_0$ for this fixed application, even though
the trace power is not known in advance.  Independent
traces in the full initial field will be counted inside the span of
actual output initials.  The degree and value bounds determine the
possible filtered output spaces; the span inclusion transfers a lower
dimension bound to them.  We state the
relative nonvanishing input first, then prove the lifting properties
before returning to its independent proof.

\subsection{A trace supplied by a strict divisorial budget}

If $J$ is a nonempty finite set of elements of a field, $J^a$ denotes the
set of products of $a$ members, with repetitions, and $\min P(J)$ denotes
$\min_{F\in J}P(F)$.

\begin{lemma}[Strict budget and trace]\label{c:trace-existence}
Let $L$ be a finitely generated field over an algebraically closed field $k$
of characteristic zero, let
$\vartheta=\nu/f^{1/n}$ be a fractional rational volume form, and let
$I\subset L^*$ be a nonempty finite system whose ratios generate $L/k$.
Let $k_I$ be a positive integer.  Suppose that every prime divisor $P$ over
a proper model of $L/k$ satisfies
\begin{equation}\label{eq:E6}
 A_\vartheta(P)>\frac{\min P(I)}{k_I}.
\end{equation}
After spreading these data, in sufficiently general reductions of arbitrarily
large characteristic $\mathfrak p\equiv1\pmod{nk_I}$ there are $l\ge1$ and
$H\in\operatorname{span}(I^{(P_0-1)/k_I})$, where $P_0=\mathfrak p^l$, such that
\begin{equation}\label{eq:E7}
 \mathcal C^l_\vartheta(H)
 :=\frac{\mathcal C^l\bigl(Hf^{(P_0-1)/n}\nu\bigr)}{\nu}\ne0.
\end{equation}
Here $\mathcal C^l$ is iterated Cartier trace on rational top forms over
perfect constants.  One may choose $H$ to be a single product from
$I^{(P_0-1)/k_I}$.
\end{lemma}

The operator in \eqref{eq:E7} is independent of the presentation of
$\vartheta$.  Indeed replacing $(\nu,f)$ by $(a\nu,a^nf)$ multiplies the
trace input by $a^{P_0}$ and the denominator by $a$, and Cartier trace is
$\mathfrak p^{-l}$-linear.  We will repeatedly use its logarithmic-coordinate formula.
If $x=(x_1,\ldots,x_s)$ is a $\mathfrak p$-basis, then
\[
 b=\sum_{0\le\alpha_j<P_0}b_\alpha^{P_0}x^\alpha
 \quad\Longrightarrow\quad
 \mathcal C^l\left(b\bigwedge_jd\log x_j\right)
     =b_0\bigwedge_jd\log x_j.
\]
This is the iterated Cartier isomorphism
\cite[Th\'eor\`eme~1]{Cartier1957}: logarithmic differentials are
fixed, and all the other displayed monomials are killed.

The proof of Lemma~\ref{c:trace-existence} appears in the final
subsection, after the lifting and finite-space estimates.  Those steps
explain how its surviving product supplies
the actual degree, value, and dimension bounds just described.


\subsection{Survival of initials and the degree bound}

We now distinguish the actual field $K$ from the full initial field $K_w$
of a real monomial valuation $w$.  The latter is the rational function
field in the independent normal variables over the field of the stratum.
We use the angular systems $V_m$ and the fractional form $\theta$ from
\eqref{eq:A13}.  The reductions in the following statements retain the
monomial chart, its stratum, its form identities, and separating coordinates.
The simultaneous finite-data construction is given below in
Lemma~\ref{c:trace-finite-data}.

\begin{lemma}[Initial survival]\label{c:trace-initial-survival}
Transport $w$ to a reduction of its monomial chart, using the same positive
normal weights.  Let $x$ consist of its normal parameters and a unit
separating tangential basis.  Retain the conditions that $x$ is a $\mathfrak p$-basis
of the actual field and that its initials form a $\mathfrak p$-basis of $K_w$.
Write
\[
 \theta=\frac{\bigwedge_jd\log x_j}{f_x^{1/n}}.
\]
For $\mathfrak p\equiv1\pmod n$, $P_0=\mathfrak p^l$, and every $G$ in the reduced actual
field,
\begin{equation}\label{eq:E8}
 w\bigl(\mathcal C^l_\theta(G)\bigr)
 \ge \frac{w(G)}{P_0}
      -\left(1-\frac1{P_0}\right)A_\theta(w).
\end{equation}
We set $w(0)=+\infty$.  For $G\ne0$, equality holds if and only if
$\mathcal C^l_{\theta_w}(\bar G)\ne0$; in that case
\[
 \overline{\mathcal C^l_\theta(G)}
       =\mathcal C^l_{\theta_w}(\bar G).
\]
The discrepancy in \eqref{eq:E8} is the characteristic-zero numerical
quantity transported by the chart formula.  The same assertions hold
with ordered monomial weights, and the inequality holds under local real
weight variation within the same open stratum cone.
\end{lemma}

\begin{proof}
The $P_0$-basis expansion has the unique form
\[
 Gf_x^{(P_0-1)/n}
   =\sum_{0\le\alpha_j<P_0}c_\alpha^{P_0}x^\alpha,
 \qquad \mathcal C^l_\theta(G)=c_0.
\]
There is no cancellation between the least-valued terms on the right:
the monomials $\bar x^\alpha$ are linearly independent over $K_w^{P_0}$.
Consequently
\[
 P_0w(c_0)\ge w(G)+\frac{P_0-1}{n}w(f_x).
\]
The logarithmic form identity \eqref{eq:A6} says
$A_\theta(w)=-w(f_x)/n$, proving \eqref{eq:E8}.
Taking least terms in the same expansion computes
$\bar G\bar f_x^{(P_0-1)/n}$ in the initial $P_0$-basis.  Its zero-index
coefficient is nonzero exactly when the inequality is an equality, and
then it is the initial of $c_0$.  The leading-form identity in
\eqref{eq:A6} identifies the fractional form in this calculation with
$\theta_w$.  This proves the equality criterion as well as the asserted
initial identity.  The argument only uses the order, addition of values,
and independence over powers, so it also applies to ordered value groups.
For nearby positive weights in the same chart the coefficient argument
again gives the inequality, with the corresponding numerical discrepancy.
\end{proof}

\begin{lemma}[Angular degree of a traced output]\label{c:trace-degree}
In the retained reductions of the angular algebra, if $G\in V_M$, then
$\mathcal C^l_\theta(G)\in V_{m'}$ for some $m'\in n\Z_{\ge0}$ with
\[
 m'\le\frac{M}{P_0}+n.
\]
\end{lemma}

\begin{proof}
Choose $m'\in n\Z_{\ge0}$ with
$P_0m'+P_0-1\ge M$ and $m'\le M/P_0+n$.  Recall the cone notation
\[
 \Omega=t\,\omega\wedge d\log t,\qquad
 h=t^nf_h^W,\qquad \theta=\omega/(f_h^W)^{1/n}.
\]
Since $\mathfrak p\equiv1\pmod n$, the exponent in
\[
 b=G h^{(P_0m'+P_0-1)/n}
\]
is integral; $b$ is regular on $W$ by $G\in V_M$ and the displayed
inequality.  Apply $b\mapsto\mathcal C^l(b\Omega)/\Omega$.
The log-$t$ formula and Frobenius-linearity give
\[
 \frac{\mathcal C^l(b\Omega)}{\Omega}
       =h^{m'/n}\mathcal C^l_\theta(G).
\]
Indeed the $t$ exponent in $b\Omega$ is $P_0(m'+1)$, which trace divides
by $P_0$, and the remaining coefficient is precisely the defining
coefficient of $\mathcal C^l_\theta$.
The output is regular in codimension one, because $\Omega$ is a regular
canonical generator there and Cartier trace sends regular top forms to
regular top forms.  It is therefore regular everywhere by normality.
Its physical degree is $m'$, so division by $h^{m'/n}$ places the output
in $V_{m'}$.  The normal graded algebra of $W$, its field and its
codimension-one canonical generator are retained in the spread.  This
argument requires no new discrepancy assertion in positive characteristic.
\end{proof}

\subsection{From a relative trace to the full initial field}

Fix independent homogeneous probes $g=(g_1,\ldots,g_{r_g})$ in $K_w$ and
use the relative fractional form obtained by dividing $\theta_w$ by the
wedge of the \emph{ordinary} differentials $dg_j$.  To describe its
geometric generic field without confusing constants and functions, first
introduce independent constants $\xi_1,\ldots,\xi_{r_g}$ and put
$k'=\overline{k(\xi_1,\ldots,\xi_{r_g})}$.  Extend the full field model by
these constants; the functions $g_j$ remain independent over $k'$ on this
total model.  At the generic point $\eta$ of a chosen component of
$g=\xi$, the residue field is the chosen geometric generic field for
$K_w/k(g)$.  Every nonzero rational function from $K_w$ before this
constant extension is a unit at $\eta$, and its residue is its prescribed
relative image, since $\eta$ maps to the original generic point.

Choose a smooth model on which the map $g$ is smooth at $\eta$.  Then
\[
 u_j'=g_j-\xi_j\quad(1\le j\le r_g)
\]
are a regular system of parameters at $\eta$.  Choose unit supplementary
coordinates $s$ whose residues form a separating transcendence basis of
the relative field.  They may be chosen among pre-extension rational
functions.  Spread the total model, this component, its relative model,
and all these coordinate and restriction data together.  After reduction,
$u',s$ are a $\mathfrak p$-basis of the total field, and the Taylor initials of these
coordinates form a $\mathfrak p$-basis of the Taylor initial field.  The relative
trace below is taken in the field of this transported component.  Its
existence is supplied by Lemma~\ref{c:trace-existence} applied before
reduction to the fixed characteristic-zero pair.

\begin{lemma}[Taylor flags]\label{c:trace-taylor-flags}
Suppose a product of restrictions of homogeneous initials gives a
nonzero relative trace as in \eqref{eq:E7}, and let $H$ denote the lift
of that product by the same factors to the full initial field.  Then
$\mathcal C^l_{\theta_w}$ is nonzero on
\begin{equation}\label{eq:E9}
 H\prod_{j=1}^{r_g}(u_j')^{P_0-1}.
\end{equation}
More precisely, give $u_j'$ the standard positive basis values in
lexicographically ordered $\Z^{r_g}$, and use the resulting Taylor valuation
at $\eta$.  If a shift input $B$ has Taylor order $P_0a$, with
$a\in\Z^{r_g}$, and has a nonzero scalar leading coefficient, then
\[
 \operatorname{ord}_{\mathrm T}\left(
 \mathcal C^l_{\theta_w}\left(
 B H\prod_j(u_j')^{P_0-1}\right)\right)=a.
\]
For $r_g=0$ the empty product is one and the statement is the original
relative nonvanishing assertion.
\end{lemma}

\begin{proof}
Write $\bar u_j'$ for the normal variables of the Taylor initial field;
this bar refers to Taylor initial, rather than the original $w$-initial.
In the logarithmic basis for $u',s$, write the denominator of $\theta_w$
as $f_{u',s}$.  Its Taylor initial is
\[
 \operatorname{in}_{\mathrm T}(f_{u',s})
       =\prod_j(\bar u_j')^{-n}f_{\mathrm{rel}},
\]
where $f_{\mathrm{rel}}$ is the denominator of
$\theta_{w,\mathrm{rel}}$ in the logarithmic basis of the residue
coordinates $s$.  This is precisely the effect of first dividing by
$\bigwedge_jdg_j=\bigwedge_jdu_j'$ and only then changing from ordinary
to logarithmic normal differentials.  The unit coefficient identities
are among the retained data at the fiber.

The factors of $H$ are units at $\eta$ and restrict to the factors of
the nonzero relative trace.  If the scalar leading coefficient of $B$
is $c\ne0$, the Taylor initial of the coefficient to which Cartier trace
is applied is consequently
\[
 c(\bar u')^{P_0a}
       H|_\eta\, f_{\mathrm{rel}}^{(P_0-1)/n}.
\]
The normal factor $\prod_j(\bar u_j')^{-(P_0-1)}$ from the denominator
has been canceled exactly by the flags in \eqref{eq:E9}.  Extraction in
the initial $P_0$-basis therefore gives
\[
 c^{1/P_0}(\bar u')^a
       \mathcal C^l_{\theta_{w,\mathrm{rel}}}(H|_\eta),
\]
which is nonzero.  The ordered-weight version of
Lemma~\ref{c:trace-initial-survival} identifies it with the Taylor initial
of the full trace, proving the asserted order and nonvanishing.
\end{proof}

There is one more lifting step.  Expand an input in
Lemma~\ref{c:trace-taylor-flags}, including any shift, as a linear
combination of products of homogeneous angular initials over the enlarged
constants.  Suppose that the factors in this expansion have actual lifts
with the prescribed degree and value bounds.  A homogeneous summand with
nonzero trace \emph{in the full initial field} then has an actual lift
$G$ for which Lemma~\ref{c:trace-initial-survival} gives equality in
\eqref{eq:E8} and identifies its traced initial.  Additivity and perfectness
of the constants show that the full initial-field outputs of all these
inputs belong to the span of such actual output initials.  In particular,
a family of independent shifted outputs is spanned in this way; there
is no need to select a single actual summand equal to each full output.
Degree and value budgets are checked for every expanded summand.  Thus
cancellation in the expansion is not needed for the bounds.

The order of the argument is essential: relative survival first gives
full initial-field survival by Lemma~\ref{c:trace-taylor-flags}, and only
then does Lemma~\ref{c:trace-initial-survival} give actual survival and
its equality formula.  Mere nonvanishing of an actual trace would not
give that equality.

Figure~\ref{fig:trace-survival} summarizes these two lifting steps. The
output degree is controlled term by term, while independent full-field
outputs are counted inside the span of actual output initials.
\begin{figure}[tbp]
\centering
\begin{minipage}{0.94\textwidth}
\centering
\fbox{\parbox{0.94\linewidth}{\centering
  \textbf{Geometric generic fibre of $K_w/k(g)$}\\[2pt]
  A strict divisorial budget gives a product $H|_\eta$ with
  nonzero relative Cartier trace.}}
\par\smallskip
$\Big\downarrow$\quad
\parbox{0.76\linewidth}{Taylor flags
  $\prod_j(g_j-\xi_j)^{P_0-1}$\\
  Lemma~\ref{c:trace-taylor-flags}}
\par\smallskip
\fbox{\parbox{0.94\linewidth}{\centering
  \textbf{Full initial field $K_w$}\\[2pt]
  The flagged product has nonzero trace. Expand it in homogeneous
  angular factors and retain the summands whose traces survive.
  Independent full-field outputs lie in the span of the resulting
  actual output initials.}}
\par\smallskip
$\Big\downarrow$\quad
\parbox{0.76\linewidth}{Lift the same factors to the original field\\
  Lemmas~\ref{c:trace-initial-survival} and~\ref{c:trace-degree}}
\par\smallskip
\fbox{\parbox{0.94\linewidth}{\centering
  \textbf{Actual function field $K$}\\[2pt]
  Each surviving summand gives an output $F'$ attaining equality
  in~\eqref{eq:E8}, with controlled physical degree.}}
\end{minipage}
\caption{The order of trace lifting. The downward steps use the retained
reductions over the extended perfect constants. Whole output spaces and their
filtration bounds are fixed in characteristic zero before the
characteristic and the Frobenius power are chosen.}
\label{fig:trace-survival}
\end{figure}


\subsection{Finite data retained before choosing the characteristic}

The trace power and the surviving products are chosen only after
reduction.  Their bounds must therefore hold for every possible output
in the prescribed degrees.  We retain the relevant finite-cutoff
filtrations of the whole modules $V_m$, including their high-cost
parts, before making those choices.

\begin{lemma}[Simultaneous finite-data reduction]\label{c:trace-finite-data}
For each single application of the preceding trace constructions one may
retain simultaneously the following data on a sufficiently general spread:
the specified actual and initial monomial charts and fractional forms;
the total and relative charts for the geometric generic component;
finitely many functions, field identities and initial identifications;
finitely many Taylor leading terms; and all specified finite-degree,
finite-cutoff comparisons between the angular filtrations of $v$ and $w$.
The last assertion concerns every element of each tested space, including
elements with arbitrary coefficients in the extended reduced constants.
A strict lower $w$-bound on the part of high $v$-cost also persists locally
under positive weight variation in the specified chart.

These data, and all required output degrees and cutoffs, are fixed before
$\mathfrak p$ and $P_0$.  They may depend on the sequence member, the optimizer,
the parameter, and the particular application.  Products formed later
from the retained finite lists, with multiplicities depending on $P_0$,
require no new spreading assertion.
\end{lemma}

\begin{proof}
\emph{Models, coordinates, and forms.}
Descend the finite algebraic data to a finitely generated integral base
over $\Z$.  When geometric generic constants $k'$ are used, first extend
the actual and full initial models by these constants as well, and include
the finitely many constants occurring in the constructions in the spread.
Spread generators, relations, maps, their inverse identities on birational
opens, and rational forms.  Shrinking the base retains flatness, geometric
integrality, and normality where used.  It also retains the smooth charts,
chosen transverse integral strata, normal parameters and \'etale differential
or supplementary coordinates.  On $W$ retain the graded algebra, its
birational field description, and the fact that $\Omega$ generates in
codimension one: its generating smooth open has complement of codimension
at least two, and the same remains true after shrinking.  Thus the
$\mathfrak p$-bases and their initial-field identifications used above coexist with
the angular algebra.  They may also be imposed together with the finite
spreading conditions in the deferred proof of
Lemma~\ref{c:trace-existence}, by a common shrinking of the spread.

\emph{Finite initial expressions.}
Consider a regular numerator or denominator at one of the stratum points.
For sufficiently large ordinary normal degree $N$, successive jets express
it modulo the $N$-th power of the normal ideal as a finite sum of distinct
normal monomials with lifted unit coefficients.  Retain their nonzero
residue coefficients and the congruence, whose remainder is a combination
of degree-$N$ monomials with regular coefficients.  Near fixed positive
normal weights, all terms in that remainder have weight greater than the
weights under consideration.  Consequently the finite expansion determines
the value and the initial expression there.  Spread this finite expansion
and its unit and stratum data.  Its minimum and homogeneous initial
expression then agree in characteristic zero and in the identified reduced
stratum/normal-variable field.  Applying this to numerators and denominators
handles the specified rational functions.  For a finite list it also
retains the piecewise leading rules for locally varying positive weights.

Finite initial linear independences are retained by a nonzero evaluation
determinant, after adjoining the finitely many points or constants needed
to witness it.  Finite algebraic independences in characteristic zero are
retained by nonzero differential wedges.  Include the full initial-field
model itself through its stratum/normal-variable identification.  This
ensures that preservation of the initials is a statement in the full
initial field, rather than only a list of numerical equalities.

\emph{The geometric generic component and Taylor data.}
Spread the smooth component chart at $g=\xi$ as a subvariety of the total
initial-field model over the enlarged constants.  Retain geometric
integrality of the selected component, the field identifications, the
local normal equations $u'=g-\xi$, the differential conditions, and the
restrictions of the specified functions and forms to that component.
For a specified Taylor leading exponent, the monomial ideal of exponents
strictly above it is finitely generated.  The required leading expansion
is a finite congruence modulo this ideal.  A congruence first obtained in
the completion contracts to the local ring, by faithful flatness of
completion.  Thus it too can be spread, together with any required nonzero
scalar leading coefficient.  Apply this to a numerator and a denominator
if necessary.  Later shift products only multiply these predetermined
leading expressions.  Their possibly $P_0$-dependent multiplicities cause
no new condition on the spread.  This retains the relative form, the
Taylor flag formula, and the actual-to-initial identification simultaneously.

\emph{Entire spaces at finite cutoffs.}
The angular space $V_m$ is a finite module over $k[Z]$.  At each of the
finitely many tested degrees, choose module generators $e_{m,t}$ and
choose affine algebra generators $a_1,\ldots,a_s$ of $k[Z]$ vanishing at
$z$.  Both $v$ and $w$ are positive on these $a_j$, by the given barriers
and the initial construction.  Retain the module generation and these
values after constant extension and reduction.  There is a common
$\delta>0$ that is a lower bound for all their $v$- and $w$-values;
for $w$ it remains a lower bound, after a harmless decrease, on a small
neighborhood of the chosen positive weights.  If
$\mathfrak a=(a_1,\ldots,a_s)$, every element of
$\mathfrak a^NV_m$ has both values at least
\[
 N\delta+\min_t\{v(e_{m,t}),w(e_{m,t})\}.
\]
Choose $N$ so this exceeds every relevant upper cutoff.  The same estimate
holds after reduction, and locally for $w$.  Modulo this tail, the finite
span of $a^\alpha e_{m,t}$ with $|\alpha|<N$ represents all of $V_m$.

Choose bases of that finite span adapted separately to the two valuation
filtrations.  Retain the relations, change-of-basis matrices, values, and
linear independence of initials in each same-value slice.  The last
condition prevents cancellation of a least-valued slice for any nonzero
coefficient vector over the extended constants.  Therefore the retained
finite spans have exactly the same filtration dimensions, gaps, and
specified cutoff inclusions.  Combined with the tail estimate, these
statements hold for every element of $V_m$ through the specified cutoffs.
For the high-$v$-cost subspace use its retained finite basis.  A lower
$w$-bound with strict margin persists under a small change of positive
weights by continuity of each basis value; the tail remains above the
same bound.  This proves also the local assertion.  Only these finite
comparisons are transported; no global optimization or infinite ratio
test is asserted in the reduction.
\end{proof}

\subsection{Applying the strict budget and accounting for all factors}

\begin{proposition}[Budgeted actual outputs]\label{c:trace-budget}
Apply the strict budget \eqref{eq:E5} on the geometric generic field of
$K_w$ over the independent probes $g$, choosing the optional $\bar U_j$
to include generators of the relative field.  Assume each probe has a fixed
angular lift $g_j=\overline{G_j}$ with $G_j\in V_{d_j}$.  Any prescribed
shift is assumed to satisfy the Taylor leading-term hypothesis of
Lemma~\ref{c:trace-taylor-flags} and to have an expansion in fixed homogeneous
factors with angular lifts and bounded factor counts divided by $P_0$.
Clear all the rational exponents
by one positive integer $k_I$.  The trace product may be chosen so that
the multiplicities of the active factors in its lift $H$ are
\begin{equation}\label{eq:E10}
 (P_0-1)\tau_i\alpha_i,\qquad
 \tau_i\ge0,\qquad
 \sum_i\tau_i=1
 \quad\text{or}\quad \sum_i\tau_i\le1\ \text{if the cap is active}.
\end{equation}
The multiplicity of a tiny factor $\bar U_j$ is at most
$(P_0-1)e_j$.  Expanding the Taylor flags adds between zero and $P_0-1$
copies of each probe.  Every surviving homogeneous initial summand has
an actual traced output $F'$ of degree at most its total physical input
degree divided by $P_0$, plus $n$.  At the chosen $w$ its value satisfies
equality in \eqref{eq:E8}; on a local affine continuation in the same
open stratum cone the lower inequality remains available.

Common mandatory factors included in every generator of the input
system remain in every product and every expanded summand.  All output
spaces and finite cutoff comparisons needed for this conclusion may be
chosen before $\mathfrak p$ and $P_0$, from bounds for the normalized factor counts.
\end{proposition}

\begin{proof}
Choose $k_I$ so that $k_I\alpha_i$ and $k_Ie_j$ are positive integers.
On the relative field take the system consisting of
\[
 \bar F_i^{k_I\alpha_i}\prod_j\bar U_j^{l_j},
 \qquad 0\le l_j\le k_Ie_j,
\]
for each active index $i$; if the cap is active include also the same
products with initial active factor $1$.  The minimum of a valuation on
this system, divided by $k_I$, is exactly
\[
 \min_{i\ \mathrm{active}}\alpha_iP(\bar F_i)
        +\sum_je_j\min\{0,P(\bar U_j)\},
\]
with zero included in the minimum for the active cap.  Its ratios include
each $\bar U_j$, by comparing the exponents $0$ and $1$ while holding
the other choices fixed.  The field-generation provision in
\eqref{eq:E5} therefore makes the system birational.  Equation
\eqref{eq:E5} is exactly the hypothesis \eqref{eq:E6} for this fixed
system.

Apply Lemma~\ref{c:trace-existence} and choose a surviving product of
$N_I=(P_0-1)/k_I$ generators.  If $N_i$ occurrences use the active index
$i$, put $\tau_i=N_i/N_I$.  Then its $\bar F_i$ exponent is
$N_ik_I\alpha_i=(P_0-1)\tau_i\alpha_i$, proving \eqref{eq:E10}.
A cap occurrence accounts for the possible deficit in $\sum_i\tau_i$.
The bound for the $\bar U_j$ exponents follows from
$N_Ik_Ie_j=(P_0-1)e_j$.

Lift this product using the same full initial factors and multiply by
the Taylor flags of \eqref{eq:E9}, and by any shift prescribed in the
application.  Each flag expands over the extended constants as a linear
combination of $g_j^a$ with $0\le a\le P_0-1$.  More generally, expand
every prescribed shift in its fixed list of homogeneous factors, with
the Taylor leading term required by Lemma~\ref{c:trace-taylor-flags}.
Lemma~\ref{c:trace-taylor-flags} supplies full initial-field survival.
For every homogeneous summand whose initial trace survives,
Lemma~\ref{c:trace-initial-survival} produces the actual trace and its
initial with equality in \eqref{eq:E8}.  Multiplying actual lifts puts
the input in $V_M$, where $M$ is the sum of their physical degrees, and
Lemma~\ref{c:trace-degree} gives $m'\le M/P_0+n$.
These surviving initials span the full initial-field output, as explained
after Lemma~\ref{c:trace-taylor-flags}.

If, in a later application, every system generator is multiplied by
$\prod_a\bar B_a^{k_I\beta_a}$, with nonnegative rational $\beta_a$
and integral $k_I\beta_a$, then every selected product contains
$\prod_a\bar B_a^{(P_0-1)\beta_a}$.  These factors also remain in each
term after expanding flags and shifts.  Selection of a surviving term
therefore retains all mandatory factors, including their degree and
value contributions.

Finally fix the finite lists of factors and their actual lifts before
reduction.  Their physical degrees, values at the tested valuations and,
when needed, Lipschitz bounds on the given continuation are now known.
For any additional shifts, include their separately verified bounds on
normalized factor counts among these data.  Each multiplicity is then
bounded by a fixed constant times $P_0$.  Summing and dividing by $P_0$,
the degree estimate and equality at $w$ give predetermined output degree
and $w$-value bounds; the rounding term $n$ is also fixed.  The required
$v$-filtration cutoffs are chosen before reduction as well.  The same
accounting is applied separately to every term in
the expansions of flags and shifts.  Any output spaces, their
$v$-filtration cutoffs, and high-$v$-to-high-$w$ inclusions used later
are retained in their entirety by Lemma~\ref{c:trace-finite-data}, before
$\mathfrak p$ and $P_0$ are chosen.  Thus they cover every possible later output
coefficient vector, without having to lift that particular vector to
characteristic zero.

The normalized contributions of the tiny factors are bounded by the
fixed finite data times $e_j$.  Those from replacing $\kappa/b_i$ by
$\alpha_i$ are bounded by the same finite data times the approximation
errors, since the $\tau_i$ lie in the simplex of \eqref{eq:E10}.
As in the construction of \eqref{eq:E5}, first choose the positive
$e_j$ sufficiently small and then approximate sufficiently closely to
respect any prescribed finite tolerances.  This controls physical degrees,
values and Lipschitz changes on the already specified continuation.
The strict budget is kept during these choices.  Neither uniform
comparison for varying pairs nor a global characteristic-$\mathfrak p$ optimizer
is used.
\end{proof}

\subsection{Proof of relative trace existence}

\begin{proof}[Proof of Lemma~\ref{c:trace-existence}]
\emph{The cone and its sub-klt boundary.}
Introduce an indeterminate $u$, put $S=k[Iu]$, and let $R_c$ be the
normalization of $S$ in its fraction field.  The ratios of $I$ generate $L$,
so $\Frac(S)=L(u)$.  Finite normalization gives a nonnegatively graded
normal affine domain $R_c$ with $(R_c)_0=k$.  The conductor is a nonzero
homogeneous ideal; multiplying a nonzero homogeneous conductor element by
a positive-degree element if necessary, choose $d_c$ of positive degree
with
\[
 0\ne d_c\in R_c,\qquad d_cR_c\subset S.
\]
Set
\[
 \Theta=\nu\wedge d\log u,\qquad
 K^\Theta=\Div_{R_c}(\Theta),\qquad
 \Delta_c=-K^\Theta+\frac1n\Div(f)-\frac1{k_I}\Div(u).
\]
In particular $K^\Theta+\Delta_c$ is rational-principal.

We verify that $(\Spec R_c,\Delta_c)$ is sub-klt.  An invariant prime
valuation $D$ centered on $\Spec R_c$ is Gaussian in $u$ over $L$.
Its restriction to $L$ is either trivial or a positive multiple of a prime
divisorial valuation $P$.  For the latter assertion, the restriction has
rational rank one, and the residue-transcendence-degree inequality for
$L(u)/L$, together with Abhyankar equality for $D$, forces Abhyankar
equality for the restriction.  We allow the positive multiple in the
notation $P$; \eqref{eq:E6} is homogeneous and therefore applies to it.
The logarithmic differential formula \eqref{eq:A6} gives
\[
 A_{R_c,\Delta_c}(D)=A_\vartheta(P)+\frac{D(u)}{k_I}.
\]
Since $D(Fu)\ge0$ for every $F\in I$, we have
$D(u)\ge-\min P(I)$, and the displayed discrepancy is positive by
\eqref{eq:E6}.  If the restriction is trivial, set $A_\vartheta(0)=0$.
Then $D(u)\ge0$, and equality would make the Gaussian valuation $D$
trivial; hence again the discrepancy is positive.  These tests suffice:
an equivariant log resolution of the homogeneous divisor data has invariant
divisor components, because a connected torus cannot permute a finite set
nontrivially.  All the divisors on this resolution consequently have
positive log discrepancy, which is the sub-klt criterion.

\emph{An effective boundary and multiplier-ideal membership.}
Choose a homogeneous $0\ne h_c\in R_c$ such that
\[
 \Delta'=\Delta_c+\Div(h_c)\ge0.
\]
We record a construction that also makes every coefficient of $\Delta'$
at most one.  Each coefficient of $\Delta_c$ on $\Spec R_c$ is strictly
less than one.  If they are all nonnegative, take $h_c=1$.  In particular,
when $L=k$, the ring is $k[u]$ and the sole boundary coefficient is
$1-1/k_I$.  Otherwise $\dim R_c\ge2$.  At the finitely many negative
prime components $Q_i$ write
$a_i=\lceil-\operatorname{coeff}_{Q_i}\Delta_c\rceil$, and use the
homogeneous ideal
\[
 \mathfrak a=\bigcap_i Q_i^{(a_i)}.
\]
Its minimum order is $a_i$ at $Q_i$ and zero at every other height-one
prime, as follows by localization.  Choose homogeneous generators of
$\mathfrak a$.  Multiply each by all products of the degree-one generators
$Iu$ that bring it to one sufficiently large common degree.  The resulting
finite system retains the same minimum order at every height-one prime.
Indeed no such prime contains all of $Iu$: by the finite extension from
$S$, the inverse image of the vertex has codimension $\dim R_c\ge2$.
A general combination attains the prescribed minimum orders at the finite
support of $\Delta_c$.  Away from that support the system has no fixed
prime divisor, and Bertini on the smooth basepoint-free open makes any
new divisorial zeros reduced.  This gives the asserted $h_c$ and boundary.

Write $\mathcal J$ for the multiplier ideal, allowing a fractional ideal
for a non-effective sub-pair.  Sub-klt gives
\[
 R_c\subseteq\mathcal J(R_c,\Delta_c).
\]
For example, this containment follows on a log resolution by rounding up
the relative discrepancy divisor, whose coefficients are greater than
$-1$.  Cartier translation therefore yields
\[
 \mathcal J(R_c,\Delta')
    =h_c\mathcal J(R_c,\Delta_c)\supseteq h_cR_c.
\]
Only the membership of $h_c$ is needed.  The fractional multiplier ideal
of a non-effective sub-klt boundary need not equal the structure sheaf.

\emph{Comparison and the full family of Cartier maps.}
Apply the fixed-pair multiplier/test-ideal comparison to the effective pair
$(\Spec R_c,\Delta')$.  Its hypotheses hold because
\[
 K^\Theta+\Delta'
   =\frac1n\Div(f)-\frac1{k_I}\Div(u)+\Div(h_c)
\]
is rational-principal, with index dividing $nk_I$; neither klt of
$\Delta'$ nor an independent $\Q$-Gorenstein hypothesis on $R_c$ is
required.  The comparison is \cite[Theorem~3.2]{Takagi}; the dense-open
form with the big test ideal is recalled in
\cite[Theorem~11]{DFDTT}, applied with
$Z=0$.  Thus, on a sufficiently general reduction to a finite residue
field of characteristic $\mathfrak p\equiv1\pmod{nk_I}$, the image of $h_c$ belongs
to the big test ideal $\tau_b(R_c,\Delta')$.

We use the following characterization: on an $F$-finite normal domain the
big test ideal of an effective rational divisor pair is the smallest
nonzero ideal preserved by all restricted maps
\[
 \Hom_{R_c}\!\left(
 F_*^eR_c\bigl(\lceil(\mathfrak p^e-1)\Delta'\rceil\bigr),R_c\right),
 \qquad e\ge1.
\]
This is \cite[Definitions~3.15--3.16 and Theorem~3.18]{SchwedeNonQG},
with the divisor-pair identification in its Remarks~3.10 and~3.17.
The log-$\Q$-Cartier comparison with the usual test ideal is also recorded
in \cite[Theorem~2.8(2)]{Takagi} and
\cite[Remark~6.1]{SchwedeNonQG}.

In the reduction define
\[
 T^e(b)=\frac{\mathcal C^e(b\Theta)}{\Theta},\qquad
 L_e=R_c\bigl((1-\mathfrak p^e)(K^\Theta+\Delta')\bigr).
\]
We use the convention
$R_c(D)=\{b\in\Frac(R_c):\Div(b)+D\ge0\}$, so explicitly
\[
 L_e=f^{(\mathfrak p^e-1)/n}u^{-(\mathfrak p^e-1)/k_I}h_c^{\mathfrak p^e-1}R_c.
\]
All the allowed maps of degree $e$ are exactly $T^e(r\,\cdot)$,
$r\in L_e$.  Here every $(\mathfrak p^e-1)\Delta'$ is integral, because
$\mathfrak p-1$ is divisible by $nk_I$; thus this statement describes the full
allowed family in every Frobenius degree, with no rounding or omitted
degrees.  This is the divisor--Frobenius-map correspondence
\cite[Corollary~3.10 and Theorems~3.11, 3.13]{SchwedeFAdjunction}.
We check it in the chosen trivialization to fix the normalization.
On the fraction field a nonzero
Frobenius-linear functional generates all such functionals by input
multiplication.  At a smooth codimension-one point, Cartier trace for a
regular canonical generator in ordinary coordinates extracts the monomial
with all exponents $\mathfrak p^e-1$.  Multiplication and trace give the perfect
pairing on the Frobenius basis, also after an \'etale change of coordinates.
Changing the canonical generator to $\Theta$ twists the multiplier by
the $(1-\mathfrak p^e)$-th power of its coefficient.  Allowing the poles prescribed
by $(\mathfrak p^e-1)\Delta'$ gives precisely $L_e$ at that height-one point.
Normality makes these codimension-one conditions sufficient globally.

\emph{A compatible ideal and homogeneous extraction.}
Set $T^0=\id$, $L_0=R_c$, and
\[
 J=\sum_{l\ge0}T^l(d_ch_cL_l)\subset R_c.
\]
This is an ideal: ordinary multiplication can be moved inside $T^l$ by
$aT^l(b)=T^l(a^{\mathfrak p^l}b)$.  It is nonzero, since its degree-zero Frobenius
summand is $d_ch_cR_c$.  It is preserved by every allowed map.  Indeed,
if $a_s\in L_s$ and $a_l\in L_l$, composition gives
\[
 T^s\bigl(a_sT^l(d_ch_ca_l)\bigr)
   =T^{s+l}\bigl(d_ch_ca_s^{\mathfrak p^l}a_l\bigr),\qquad
 L_s^{\mathfrak p^l}L_l\subseteq L_{s+l}.
\]
In particular the conductor factor stays present.  The minimality
characterization therefore gives $\tau_b(R_c,\Delta')\subseteq J$.
A comparison localized at the homogeneous vertex would already suffice:
compatibility localizes, because Hom localizes and denominators move
inside trace after taking Frobenius powers.

The ideal $J$ is homogeneous.  Indeed log-$u$ trace kills a homogeneous
input whose degree is not divisible by $\mathfrak p^l$, and otherwise divides its
degree by $\mathfrak p^l$.  Thus local membership of the homogeneous element $h_c$
also gives global membership: multiply by a vertex unit, and project the
result to degree $\deg h_c$; the unit's nonzero constant part is the only
contribution in that degree.  We conclude $h_c\in J$.

For $P_0=\mathfrak p^l$ the input factor $h_cL_l$ contains $h_c^{P_0}$.
Dividing the finite trace identity for $h_c\in J$ by $h_c$, using
Frobenius-linearity in every summand, gives
\[
 1\in\sum_{l\ge0}T^l\left(
  f^{(\mathfrak p^l-1)/n}u^{-(\mathfrak p^l-1)/k_I}d_cR_c\right).
\]
Project to homogeneous degree zero.  The term with $l=0$ contributes
nothing, since $d_c$ has positive degree and $R_c$ is nonnegatively
graded.  At least one term with $l\ge1$ has nonzero degree-zero output.
It arises from a homogeneous element
\[
 H u^{(P_0-1)/k_I}\in d_cR_c\subset k[Iu].
\]
Consequently $H\in\operatorname{span}(I^{(P_0-1)/k_I})$.
The log-$u$ coefficient formula identifies its nonzero output with
\eqref{eq:E7}.  Additivity and perfectness of the constants let us choose
a single product in that span with nonzero output.

\emph{The reduction choices.}
Spread the finite normalization algebra, its generators and inclusions,
its birational field identification, the conductor, the forms, and all
the boundary data together.  Shrinking the base preserves geometric
normality, grading, codimension-one orders, effectivity of $\Delta'$,
and the displayed rational-principal log canonical relation.  For the
form and divisor identities, spread their rational sections on smooth
opens, or on the chosen resolving model, and retain codimension at least
two for the complements used in codimension-one tests downstairs.
The fixed-pair comparison above then holds on a sufficiently general
open of this spread.  These requirements are compatible with any finite
additional open conditions.  A nonempty open of a finite-type integral
base of characteristic zero has nonempty fibers at all sufficiently
large primes; choose arbitrarily large primes
$\mathfrak p\equiv1\pmod{nk_I}$ and closed points in those fibers.  Their residue
fields are finite and hence $F$-finite.  After the trace identity has been
obtained, extension to algebraically closed perfect constants preserves
it, and preserves its particular nonzero output, by the coefficient
formula and faithful scalar extension.  This proves the lemma.
\end{proof}

The relative nonvanishing input is now established.  Together with the
lifting and finite-space arguments, it proves the budgeted-output
conclusions of Proposition~\ref{c:trace-budget}, which we apply next to
move the cheap probe values to zero and then to bound all angular values.
